% cuadro C12_1
\begin{table}[htbp]
\centering
\scriptsize
\caption{Anexos transversales del proyecto 2026}
\label{cua:C12_1}
\begin{tabular}{p{9.6cm}rrr}
\toprule
\textbf{Anexo} & \textbf{Millones} & \textbf{\% del etiquetado} & \textbf{\% del programable} \\
\midrule
Recursos para la atención de niñas, niños y adolescentes & 1 099 514,0 & 21,1 & 15,7 \\
Erogaciones para el Desarrollo de los Jóvenes & 735 177,8 & 14,1 & 10,5 \\
Recursos para la Atención de Grupos Vulnerables & 687 765,0 & 13,2 & 9,8 \\
Erogaciones para la Igualdad entre Mujeres y Hombres & 599 145,4 & 11,5 & 8,5 \\
Programa Especial Concurrente para el Desarrollo Rural Sustentable & 525 634,1 & 10,1 & 7,5 \\
Consolidación de una Sociedad de Cuidados & 466 674,9 & 8,9 & 6,7 \\
Acciones para la Prevención del Delito, Combate a las adicciones, Rescate de Espacios Públicos y Promoción de Proyectos Productivos & 465 729,8 & 8,9 & 6,6 \\
Erogaciones para para la implementación del Desarrollo Integral, Intercultural y Sostenible de los Pueblos y Comunidades Indígenas y Afromexicanas & 234 782,5 & 4,5 & 3,3 \\
Recursos para la Adaptación y Mitigación de los Efectos del Cambio Climático & 212 569,7 & 4,1 & 3,0 \\
Programa de Humanidades, Ciencias, Tecnologías e Innovación & 160 843,8 & 3,1 & 2,3 \\
Estrategia Nacional de Transición Energética & 17 868,4 & 0,3 & 0,3 \\
Prevención, Detección, Investigación y Sanción de Hechos de Corrupción & 9 599,1 & 0,2 & 0,1 \\
\bottomrule
\end{tabular}
\\[2pt]\begin{minipage}{0.92\textwidth}\scriptsize El total no es dinero adicional: es una suma de etiquetas, y 170 de los 288 programas etiquetados aparecen en más de un anexo. Fuente: base de anexos transversales del PPEF 2026, datos abiertos de la Secretaría. Elaborado por el ITED.\end{minipage}
\end{table}
